% !TEX root = ../Thesis.tex

\chapter{Einleitung}
\label{ch:einleitung}%
\index{Observability}%
\index{Four Golden Signals}%
\index{frag.jetzt}%

Webbasierte Anwendungssysteme werden heute häufig als verteilte, cloud-native Anwendungen umgesetzt, die aus mehreren lose gekoppelten Services bestehen. Diese Architekturform erlaubt eine unabhängige Weiterentwicklung einzelner Komponenten sowie eine flexible Skalierung, erhöht jedoch zugleich den Aufwand für Betrieb, Überwachung und Fehleranalyse im Produktivsystem \cite{faseeha2025observability,sakinala2025observability}. Insbesondere bei Microservices entstehen verteilte Fehlerzustände und komplexe Abhängigkeiten zwischen Diensten, wodurch Ursachen von Störungen nicht mehr eindeutig einzelnen Komponenten zugeordnet werden können \cite{fischer2024microservice}.

Vor diesem Hintergrund gewinnt Observability als Erweiterung klassischer Monitoring-Konzepte an Bedeutung. Ziel von Observability ist es, den internen Zustand eines Systems anhand extern erfassbarer Anzeichen nachvollziehen zu können. Dazu werden Metriken, Logs und Traces kombiniert, um Zusammenhänge zwischen Systemverhalten, Fehlern und Leistungsmerkmalen sichtbar zu machen \cite{faseeha2025observability}. Empirische Untersuchungen aus dem Bereich cloud-nativer Anwendungen zeigen, dass ein strukturierter Einsatz von Observability-Praktiken die Analyse von Incidents unterstützt und die Reaktionszeit im Betrieb verkürzen kann \cite{giamattei2024monitoring,sakinala2025observability}.

Ein häufig verwendeter Orientierungsrahmen zur Strukturierung servicebezogener Beobachtungsdaten sind die von Google im Kontext des Site Reliability Engineering (SRE) beschriebenen \emph{Four Golden Signals}: Latenz, Traffic, Fehler und Sättigung \cite{ewaschukMonitoringDistributedSystems}. Diese Signalarten erfassen grundlegende technische Symptome eines Dienstes und eignen sich als Ausgangspunkt zur Ableitung erster Monitoring-Indikatoren. Für den produktiven Betrieb liefern sie jedoch primär Rohinformationen, die erst durch Kontext, Interpretation und Bewertung für Menschen handlungsrelevant werden. Google SRE weist selbst darauf hin, dass die Wirksamkeit solcher Signale maßgeblich davon abhängt, wie sie in Alerting, Entscheidungsprozesse und betriebliche Abläufe eingebettet sind \cite{ewaschukMonitoringDistributedSystems}.

Die Open-Source-Plattform \emph{frag.jetzt} verfügt über eine solide Entwicklungs- und Deployment-Infrastruktur, jedoch ist die Erfassung und Auswertung von Betriebsdaten derzeit nicht als durchgängiges Observability-Konzept ausgelegt. Metriken, Logs, Traces und Alerting werden nicht konsistent entlang des gesamten Anfragepfads betrachtet, wodurch Zusammenhänge zwischen Nutzerverhalten, Backend-Verarbeitung und externen Abhängigkeiten nur eingeschränkt nachvollziehbar sind. Diese Einschätzung stützt sich auf die öffentlich verfügbare Projektdokumentation und Architekturübersicht des frag.jetzt-Repositories \cite{fragjetztRepo} sowie auf allgemeine Beobachtungen zu typischen Defiziten in vergleichbaren Systemen \cite{giamattei2024monitoring}.\\

Ziel dieser Arbeit ist die Entwicklung einer Observability-Strategie für \emph{frag.jetzt}, bei der technische Messgrößen systematisch mit nutzungsbezogenen Fragestellungen verknüpft werden. Die Four Golden Signals dienen dabei als unterstützender Referenzrahmen zur Strukturierung der Beobachtungsdaten, nicht jedoch als alleinige Grundlage für betriebliche Entscheidungen. Die Signale werden servicespezifisch für das PWA-Frontend, das Spring-Boot-Backend, die PostgreSQL-Datenbank sowie angebundene KI-Dienste konkretisiert und in messbare Indikatoren überführt. Ergänzend wird untersucht, welche Open-Source-Observability-Werkzeuge die Anforderungen von \emph{frag.jetzt} im Hinblick auf Erfassung und betriebliche Nutzbarkeit am besten erfüllen.

Ein weiterer Schwerpunkt der Arbeit liegt auf der Ausgestaltung eines Alerting-Ansatzes, der Auffälligkeiten innerhalb der Plattform mit Fragen der Relevanz und Dringlichkeit verknüpft. Dabei steht nicht allein die Detektion von Abweichungen im Vordergrund, sondern insbesondere die Bewertung, ob eine Handlung erforderlich ist, wen ein Problem betrifft und welche Auswirkungen es auf die Nutzung der Plattform hat. Forschungsarbeiten zeigen, dass eine Kombination aus priorisiertem Alerting und datenbasierten Analyseverfahren dazu beitragen kann, irrelevante Alarmierungen zu reduzieren und menschliche Entscheidungsprozesse im Betrieb zu unterstützen \cite{jalalvand2024alert}. Aufbauend auf den definierten Alerts werden Runbooks konzipiert, die strukturierte Diagnose- und Handlungsschritte für wiederkehrende Incident-Szenarien bereitstellen.

\section{Forschungsfragen}
\label{sec:forschungsfragen}

Die Forschungsfragen dieser Arbeit lauten:

\begin{enumerate}
    \item FF1: Wie können die Four Golden Signals für die spezifischen Services von \emph{frag.jetzt} (PWA-Frontend, Spring Boot Backend, PostgreSQL, KI-Services) konkret definiert und instrumentiert werden?
    \item FF2: Welche Observability-Stack-Kombination (Prometheus/Grafana, ELK Stack, Jaeger) ist am besten geeignet für die Monitoring-Anforderungen von \emph{frag.jetzt}?
    \item FF3: Wie kann ein intelligentes Alert-System entwickelt werden, das False Positives minimiert und actionable insights für verschiedene Stakeholder-Rollen (DevOps, Entwickler, Product Owner) bereitstellt?
\end{enumerate}
